\chapter{Appendix: Math and Data-Structure Background}
\label{ch:appendix}

% --- chapter-local shorthands (do not clash with common.tex) ---
\providecommand{\lowbit}{\operatorname{lowbit}}
\newcommand{\E}{\operatorname{\mathbb{E}}}
\newcommand{\Var}{\operatorname{Var}}
\renewcommand{\Pr}{\operatorname{\mathbb{P}}}
\providecommand{\diag}{}\renewcommand{\diag}{\operatorname{diag}}
\providecommand{\tr}{}\renewcommand{\tr}{\operatorname{tr}}
\newcommand{\amort}{\widehat{c}}
\newcommand{\find}{\textsc{Find}}
\newcommand{\union}{\textsc{Union}}

The chapters in this book rely on a small set of recurring mathematical tools: a way to talk about how fast
an algorithm runs, a handful of core data structures that show up over and over, and just enough
linear algebra and probability to follow the spectral and randomized material. None of it is
graph theory proper, which is why it lives here at the back rather than scattered through the
main text. Every tool here is used directly in the chapters ahead. Union-find powers Kruskal's
minimum-spanning-tree algorithm and almost every offline connectivity computation.
Heaps drive Dijkstra and Prim. Segment trees and Fenwick trees power the range queries that
heavy-light decomposition and Euler-tour techniques reduce to. Eigenvalues of the Laplacian
\emph{are} the content of Chapters~31--35, culminating in Chapter~41's electrical flows. And the probability section supplies what you need to follow
Karger's contraction algorithm, streaming sketches, and Chapters~35--36 (random walks and PageRank).

The treatment here is self-contained but brisk. Where a topic deserves a whole book of its
own -- the linear algebra especially -- we build the intuition and prove the facts the graph
algorithms actually use, then point to the companion \emph{Linear Algebra} volume for the
full development. Read this appendix once front to back to calibrate notation, or dip in
when a main chapter sends you here.

\chapterline

\section{Asymptotics: Big-O, Big-Theta, Big-Omega}
\label{ch42:asymptotics}

We never care about the exact running time of an algorithm down to the machine cycle. We
care how the cost grows as the input grows, and we want a notation that throws away the
clutter -- constant factors, lower-order terms, the speed of the particular computer -- and
keeps the growth rate. That is what asymptotic notation does.

\begin{intuitionbox}
Think of two algorithms as two functions $f(n)$ and $g(n)$ giving the number of basic
operations on inputs of size $n$. Big-O is a claim about the \emph{ceiling}: $f = \Oh(g)$
says ``$f$ grows no faster than $g$, ignoring constants.'' Big-Omega is the \emph{floor}:
$f=\Om(g)$ says ``$f$ grows at least as fast as $g$.'' Big-Theta is both at once: $f=\Th(g)$
pins the growth rate exactly. When we say merge sort is $\Th(n\log n)$ we mean its cost is
sandwiched between two constant multiples of $n\log n$ for all large $n$.
\end{intuitionbox}

\begin{definition}[Asymptotic classes]
\label{ch42:def-bigo}
Let $f,g:\N\to\R_{\ge 0}$. Then
\[
f=\Oh(g) \iff \exists\, c>0,\ n_0:\ \forall n\ge n_0,\ f(n)\le c\,g(n),
\]
\[
f=\Om(g) \iff \exists\, c>0,\ n_0:\ \forall n\ge n_0,\ f(n)\ge c\,g(n),
\]
\[
f=\Th(g) \iff f=\Oh(g)\ \text{and}\ f=\Om(g).
\]
We write $f=o(g)$ when $f(n)/g(n)\to 0$, a \emph{strictly} smaller growth rate.
\end{definition}

The ``$=$'' in $f=\Oh(g)$ is a historical abuse; $\Oh(g)$ is really a \emph{set} of
functions and ``$f=\Oh(g)$'' means $f\in\Oh(g)$. You can chain it ($n^2+3n=n^2+\Oh(n)$) but
never read it as an equation you can flip.

\begin{example}[Discarding the small stuff]
\label{ch42:ex-discard}
Let $f(n)=3n^2+50n+1000$. Then $f=\Th(n^2)$. Upper bound: for $n\ge 50$ we have
$50n\le n^2$ and $1000\le n^2$, so $f(n)\le 3n^2+n^2+n^2=5n^2$, giving $\Oh(n^2)$ with $c=5$,
$n_0=50$. Lower bound: $f(n)\ge 3n^2$ always, giving $\Om(n^2)$ with $c=3$. The linear and
constant terms vanish into the constant; only the leading $n^2$ survives.
\end{example}

A few facts we use constantly, all immediate from the definition:

\begin{itemize}
\item \textbf{Transitivity.} $f=\Oh(g)$ and $g=\Oh(h)$ imply $f=\Oh(h)$.
\item \textbf{Sum rule.} $\Oh(f)+\Oh(g)=\Oh(\max(f,g))$. Two phases of an algorithm
cost the larger of the two.
\item \textbf{Product rule.} If $f_1=\Oh(g_1)$ and $f_2=\Oh(g_2)$ then
$f_1 f_2=\Oh(g_1 g_2)$. A loop of $\Oh(g_1)$ iterations each costing $\Oh(g_2)$ runs in
$\Oh(g_1 g_2)$.
\item \textbf{Logs.} $\log_a n=\Th(\log_b n)$ for any constant bases, so we drop the base
inside $\Oh$. But the base \emph{matters} in an exponent: $2^n$ and $3^n$ are not
asymptotically equal.
\end{itemize}

\begin{trapbox}
For a loop running $i$ from $1$ to $n$ with variable inner cost, multiplying the iteration count
by the maximum inner cost can severely overcount. When the inner work is $\Oh(i)$, the exact sum is
$\sum_{i=1}^n \Oh(i)=\Oh(\sum i)=\Oh(n^2/2)=\Oh(n^2)$, which happens to match $\Oh(n\cdot n)$ here,
whereas inner work that drops geometrically (such as $\Oh(n/2^i)$ across tree levels) sums to $\Oh(n)$
rather than $\Oh(n\log n)$. Always sum the actual per-iteration cost instead of multiplying
the worst-case iteration by the loop count.
\end{trapbox}

\begin{connectionbox}
The notation $\Oh$ dates to Bachmann (1894) and Edmund Landau in analytic number
theory. Donald Knuth imported it into algorithm analysis in the 1960s and, in a 1976 note,
defined the matching $\Om$ and $\Th$ so that algorithm designers could state lower bounds and tight
classifications cleanly rather than abusing $\Oh$ to mean ``order exactly.'' The convention throughout
this text follows Knuth: $\Oh$ for upper bounds, $\Om$ for lower bounds, and $\Th$ when the growth
rate is pinned tightly from above and below.
\end{connectionbox}

\subsection{A hierarchy to memorize}

For large $n$, these grow in strictly increasing order, each $o$ of the next:
\[
1\ \prec\ \log\log n\ \prec\ \log n\ \prec\ n^{\eps}\ \prec\ n\ \prec\ n\log n\ \prec\ n^2\
\prec\ n^3\ \prec\ 2^n\ \prec\ n!\ \prec\ n^n,
\]
where $0<\eps<1$. Polynomials beat any polylog; exponentials beat any polynomial. The
boundary between $\poly(n)$ (tractable) and $2^{\Om(n)}$ (intractable) is the line the
complexity chapters draw.

\section{Amortized Analysis}
\label{ch42:amortized}

Some data structures are cheap on average but occasionally expensive. A dynamic array
doubles when full: most insertions are $\Oh(1)$, but the doubling insertion copies
everything. Worst-case-per-operation says $\Oh(n)$, which is true but misleading -- a
\emph{sequence} of $n$ insertions costs only $\Oh(n)$ total, so $\Oh(1)$ \emph{amortized}.
Amortized analysis is the rigorous way to charge for a sequence rather than for the worst
single step.

\begin{intuitionbox}
Imagine a prepaid coffee card. Each cheap operation overpays a little, banking credit. When
an expensive operation arrives, it spends the banked credit instead of charging you a spike.
If you can guarantee the bank never goes negative, then the total you actually paid (cheap
overpayments) upper-bounds the total cost, and dividing by the number of operations gives a
smooth per-operation price even though individual operations were lumpy.
\end{intuitionbox}

There are three standard methods; they give the same bounds but feel different.

\paragraph{Aggregate method.} Bound the total cost $T(n)$ of any sequence of $n$ operations,
then declare the amortized cost $T(n)/n$. For the doubling array, the copies happen at sizes
$1,2,4,\dots,2^{\lfloor\log n\rfloor}$, summing to $\le 2n$ copies total, so $n$ insertions
cost $\Oh(n)$ and the amortized cost is $\Oh(1)$.

\paragraph{Accounting method.} Assign each operation an \emph{amortized charge} $\amort_i$,
possibly different from its real cost $c_i$. The charges must dominate every prefix:
$\sum_{i=1}^k \amort_i \ge \sum_{i=1}^k c_i$ for all $k$. Then $\sum \amort_i$ bounds the
true total. For the doubling array, charge each insertion three credits: one for its own
write, and two banked to pay for copying itself and one earlier element when the next
doubling hits. The bank is always nonnegative, so $\amort=\Oh(1)$.

\paragraph{Potential method.} The slickest. Define a \emph{potential} function $\Phi$
mapping each data-structure state to a nonnegative real, with $\Phi(D_0)=0$. The amortized
cost of the $i$-th operation is
\[
\amort_i = c_i + \Phi(D_i) - \Phi(D_{i-1}).
\]
Summing telescopes:
\[
\sum_{i=1}^n \amort_i = \sum_{i=1}^n c_i + \Phi(D_n) - \Phi(D_0)
= \sum_{i=1}^n c_i + \Phi(D_n) \ge \sum_{i=1}^n c_i,
\]
using $\Phi(D_n)\ge 0=\Phi(D_0)$. So the amortized total bounds the real total. The hard part is
choosing $\Phi$ so that it absorbs the spikes.

\begin{recoverybox}
\textbf{Pause and think}: Why does charging 3 credits per insertion in a dynamic array guarantee that the bank balance never goes negative? After the array doubles from size $k$ to $2k$, every new insertion contributes 2 banked credits. By the time $k$ new items arrive and the array fills, exactly $2k$ credits have been saved -- precisely enough to pay for copying all $2k$ items into the next doubled array.
\end{recoverybox}

Figure~\ref{ch42:fig-amortized-bank} visualizes this credit cycle across a doubling sequence: while the true insertion cost $c_i$ spikes whenever the array resizes, the banked potential $\Phi(D_i) = 2n - \text{cap}$ climbs steadily by $+2$ on cheap insertions and drops to $0$ to pay for each copy spike, leaving a flat amortized charge $\amort_i = 3$.

\begin{figure}[H]
\centering
\begin{tikzpicture}[
  scale=0.9,
  cell/.style={draw=sourceblue!80!black,fill=sourceblue!12,minimum width=6.5mm,minimum height=5.5mm,inner sep=0pt,font=\scriptsize\bfseries},
  newcell/.style={draw=targetgreen!85!black,fill=targetgreen!18,minimum width=6.5mm,minimum height=5.5mm,inner sep=0pt,font=\scriptsize\bfseries},
  coin/.style={rectangle,rounded corners=2pt,draw=quotientorange!90!black,fill=quotientorange!22,minimum width=5.5mm,minimum height=4mm,inner sep=1pt,font=\tiny\bfseries},
  panel/.style={draw=gray!30,fill=gray!3,rounded corners=3pt},
  head/.style={font=\scriptsize\bfseries,text=sourceblue!90!black}
]
  % Panel 1: After doubling to cap=4, 2 old items + 2 new items
  \begin{scope}[shift={(0,0)}]
    \draw[panel] (-0.4,-0.85) rectangle (4.0,2.2);
    \node[head] at (1.8,1.85) {1. Fill capacity $2k=4$};
    \node[cell] (c1) at (0.4,0.6) {$x_1$};
    \node[cell] (c2) at (1.3,0.6) {$x_2$};
    \node[newcell] (c3) at (2.2,0.6) {$x_3$};
    \node[newcell] (c4) at (3.1,0.6) {$x_4$};
    % Banked credits on new items
    \node[coin] at (2.2,1.28) {$+2$};
    \node[coin] at (3.1,1.28) {$+2$};
    \node[font=\tiny,text=gray!80!black,align=center] at (1.8,-0.35)
      {$x_3,x_4$ bank $2$ credits each\\potential $\Phi = 4$ credits saved};
  \end{scope}

  \draw[-{Stealth[length=2.2mm]},very thick,gray!70] (4.2,0.65) --
    node[above=2pt,font=\tiny\bfseries]{insert $x_5$}
    node[below=2pt,font=\tiny,text=trapred]{$c_5{=}5$} (5.4,0.65);

  % Panel 2: Resize to cap=8, spend 4 banked credits + 1 write + bank 2 on x5
  \begin{scope}[shift={(6.0,0)}]
    \draw[panel] (-0.4,-0.85) rectangle (6.0,2.2);
    \node[head,text=targetgreen!85!black] at (2.8,1.85) {2. Resize to $4k=8$ ($\Phi$ pays $4$ copies)};
    \node[cell] at (0.3,0.6) {$x_1$};
    \node[cell] at (1.0,0.6) {$x_2$};
    \node[cell] at (1.7,0.6) {$x_3$};
    \node[cell] at (2.4,0.6) {$x_4$};
    \node[newcell] at (3.1,0.6) {$x_5$};
    \node[cell,fill=white,draw=gray!45,dashed] at (3.8,0.6) {};
    \node[cell,fill=white,draw=gray!45,dashed] at (4.5,0.6) {};
    \node[cell,fill=white,draw=gray!45,dashed] at (5.2,0.6) {};
    \node[coin] at (3.1,1.28) {$+2$};
    \node[font=\tiny,text=gray!80!black,align=center] at (2.8,-0.35)
      {$\Delta\Phi = 2 - 4 = -2$ absorbs copy spike:\\$\widehat{c}_5 = c_5 + \Delta\Phi = 5 - 2 = 3$};
  \end{scope}
\end{tikzpicture}
\caption{Amortized accounting and potential method on a doubling array ($\amort_i = 3$ credits per insertion). Left: after resizing from $2$ to $4$, inserting $x_3$ and $x_4$ uses $1$ credit each for the write and banks $2$ credits each ($\Phi=4$). Right: when $x_5$ triggers a resize to capacity $8$ (true cost $c_5 = 4\text{ copies} + 1\text{ write} = 5$), the $4$ banked credits pay for copying $x_1\dots x_4$, while $x_5$'s own $3$ credits pay $1$ for writing $x_5$ and bank $2$ fresh credits ($\Phi=2$).}
\label{ch42:fig-amortized-bank}
\end{figure}

\begin{example}[Binary counter]
\label{ch42:ex-counter}
Increment a $k$-bit binary counter $n$ times from zero. A single increment can flip many
bits (e.g. $0111\!\to\!1000$ flips four). Worst case per increment is $\Oh(k)$. Let $\Phi$ be
the number of $1$ bits. An increment that flips $t$ trailing ones to zero and sets one zero
to one has real cost $c_i=t+1$ and changes the bit count by $-t+1$, so
\[
\amort_i = (t+1) + (-t+1) = 2 = \Oh(1).
\]
The potential exactly cancels the cascade. Total cost of $n$ increments is $\Oh(n)$, not
$\Oh(nk)$.
\end{example}

\begin{whybox}
Why is the potential method trustworthy? Because the telescoping sum is an identity, not an
estimate. The only inequality used is $\Phi(D_n)\ge\Phi(D_0)$, which holds whenever $\Phi$ is
nonnegative and starts at zero. So \emph{any} valid potential yields a correct amortized
bound; a cleverer $\Phi$ just yields a tighter one. This is exactly the mechanism behind the
$\Oh(\log n)$ amortized bounds for splay trees and the near-constant bounds for union-find
below.
\end{whybox}

\section{Union-Find (Disjoint-Set Union)}
\label{ch42:unionfind}

Maintain a partition of $\{1,\dots,n\}$ into disjoint sets under two operations: $\find(x)$
returns a canonical representative of $x$'s set, and $\union(x,y)$ merges the two sets.
Equality of representatives answers ``are $x$ and $y$ in the same set?'' This is precisely
what Kruskal's algorithm needs to test whether an edge would close a cycle, and what offline
dynamic-connectivity and many percolation computations are built on.

\begin{intuitionbox}
Represent each set as a rooted tree whose root is the representative. $\find$ walks up to the
root. $\union$ hangs one root under the other. Two tricks keep the trees flat. \emph{Union by
rank} always hangs the shorter tree under the taller, so heights grow slowly. \emph{Path
compression} re-points every node visited during a $\find$ straight at the root, so future
finds are instant. Together they make the forest almost completely flat.
\end{intuitionbox}

\begin{figure}[H]
\centering
\begin{tikzpicture}[
  scale=0.9,
  vnode/.style={circle,draw=sourceblue!80!black,fill=sourceblue!15,minimum size=6.5mm,inner sep=0pt,font=\footnotesize\bfseries},
  rootnode/.style={circle,draw=sourceblue!90!black,fill=sourceblue!30,minimum size=7mm,inner sep=0pt,font=\footnotesize\bfseries},
  pathnode/.style={circle,draw=quotientorange!90!black,fill=quotientorange!25,minimum size=6.5mm,inner sep=0pt,font=\footnotesize\bfseries},
  header/.style={font=\footnotesize\scshape,text=sourceblue!90!black}
]
  % Panel A: Before Find(d)
  \begin{scope}[shift={(0,0)}]
    \node[header] at (1.0, 1.2) {Before $\textsc{Find}(d)$};
    \node[rootnode] (a) at (1.0, 0.4) {$a$};
    \node[pathnode] (b) at (0.2, -0.9) {$b$};
    \node[vnode] (c) at (1.8, -0.9) {$c$};
    \node[pathnode] (d) at (-0.4, -2.2) {$d$};
    \node[vnode] (e) at (0.8, -2.2) {$e$};

    % Tree pointers (child to parent)
    \draw[->,quotientorange!90!black,line width=1.6pt] (d) -- (b);
    \draw[->,quotientorange!90!black,line width=1.6pt] (b) -- (a);
    \draw[->,gray!60,thick] (c) -- (a);
    \draw[->,gray!60,thick] (e) -- (b);

    \node[font=\scriptsize,text=quotientorange!90!black] at (1.0, -2.9) {Traversed path: $d \to b \to a$};
  \end{scope}

  % Center Arrow
  \draw[->,very thick,gray!60] (3.0, -0.9) -- node[above,font=\tiny\bfseries]{Path} node[below,font=\tiny\bfseries]{Compression} (4.6, -0.9);

  % Panel B: After Find(d)
  \begin{scope}[shift={(5.6,0)}]
    \node[header] at (1.0, 1.2) {After $\textsc{Find}(d)$};
    \node[rootnode] (a2) at (1.0, 0.4) {$a$};
    \node[pathnode] (b2) at (-0.4, -0.9) {$b$};
    \node[pathnode] (d2) at (0.8, -0.9) {$d$};
    \node[vnode] (c2) at (2.2, -0.9) {$c$};
    \node[vnode] (e2) at (-0.4, -2.2) {$e$};

    % Compressed pointers directly to root
    \draw[->,targetgreen!80!black,line width=1.6pt] (d2) -- (a2);
    \draw[->,targetgreen!80!black,line width=1.6pt] (b2) -- (a2);
    \draw[->,gray!60,thick] (c2) -- (a2);
    % e was unvisited on path of Find(d), so still points to b
    \draw[->,gray!60,thick] (e2) -- (b2);

    \node[font=\scriptsize,text=targetgreen!80!black] at (1.0, -2.9) {$d$ and $b$ re-point directly to $a$};
  \end{scope}
\end{tikzpicture}
\caption{Disjoint-Set Union (DSU) path compression during $\textsc{Find}(d)$. Left: walking up from leaf $d$ to root representative $a$ visits nodes $d$ and $b$. Right: as the recursion unwinds, all traversed nodes ($d$ and $b$) re-point their parent pointers directly to root $a$, flattening the tree for all subsequent operations. The untouched branch $e$ remains attached to $b$.}
\label{ch42:fig-dsu-compression}
\end{figure}

\begin{algorithm}[H]
\caption{Union-Find with union by rank and path compression}
\DontPrintSemicolon
\SetKwFunction{FMake}{MakeSet}\SetKwFunction{FFind}{Find}\SetKwFunction{FUnion}{Union}
\SetKwProg{Fn}{function}{:}{}
\Fn{\FMake{$x$}}{$\textit{parent}[x]\gets x$;\quad $\textit{rank}[x]\gets 0$\;}
\Fn{\FFind{$x$}}{
  \If{$\textit{parent}[x]\ne x$}{$\textit{parent}[x]\gets \FFind{\textit{parent}[x]}$\tcp*{path compression}}
  \Return $\textit{parent}[x]$\;
}
\Fn{\FUnion{$x,y$}}{
  $rx\gets \FFind{x}$;\quad $ry\gets \FFind{y}$\;
  \lIf{$rx=ry$}{\Return}
  \If{$\textit{rank}[rx]<\textit{rank}[ry]$}{swap $rx\leftrightarrow ry$\tcp*{attach smaller under larger}}
  $\textit{parent}[ry]\gets rx$\;
  \lIf{$\textit{rank}[rx]=\textit{rank}[ry]$}{$\textit{rank}[rx]\gets\textit{rank}[rx]+1$}
}
\end{algorithm}

\begin{labbox}
Trace it. Start with singletons $1\dots6$, all rank $0$. Process the edge set
$\{(1,2),(3,4),(2,4),(5,6),(2,6)\}$ as Kruskal would.
\begin{itemize}
\item $\union(1,2)$: ranks tie, $2$ hangs under $1$, $\textit{rank}[1]=1$. Sets: $\{1,2\}$.
\item $\union(3,4)$: $4$ under $3$, $\textit{rank}[3]=1$. Sets: $\{1,2\},\{3,4\}$.
\item $\union(2,4)$: $\find(2)=1$, $\find(4)=3$, ranks tie ($1$ each), $3$ hangs under $1$,
$\textit{rank}[1]=2$. Set: $\{1,2,3,4\}$.
\item $\union(5,6)$: $6$ under $5$, $\textit{rank}[5]=1$. Sets: $\{1,2,3,4\},\{5,6\}$.
\item $\union(2,6)$: $\find(2)=1$ (rank 2), $\find(6)=5$ (rank 1), so $5$ hangs under $1$.
Now $\find(6)$ next time compresses $6$ straight to $1$. One set: $\{1,2,3,4,5,6\}$.
\end{itemize}
Notice how rank tracks an \emph{upper bound} on height, and how a later $\find$ flattens the
chain $6\!\to\!5\!\to\!1$ into $6\!\to\!1$.
\end{labbox}

\subsection{Why it is fast}

\begin{theorem}[Tarjan, 1975]
\label{ch42:thm-uf}
Starting from $n$ singletons, any sequence of $m\ge n$ \find\ and \union\ operations with
union by rank and path compression runs in $\Oh(m\,\alpha(n))$ time, where $\alpha$ is the
inverse Ackermann function.
\end{theorem}

The function $\alpha(n)$ grows so slowly it is at most $4$ for every $n$ that could be
written down in this universe; for practical purposes union-find is constant-time per
operation. A full proof is intricate; we prove the clean intermediate bounds and state the
sharp one.

\begin{recoverybox}
\textbf{Pause and think}: If path compression flattens find paths so aggressively, why isn't the amortized cost per operation strictly $\Oh(1)$? Because unions can link two deep trees without flattening them; an adversary can interleave unions and finds so that compressing one path leaves many long paths untouched in other branches, preventing the overall cost from collapsing to a constant.
\end{recoverybox}

\begin{lemma}[Rank bounds]
\label{ch42:lem-rank}
Under union by rank (even without path compression), a root of rank $r$ has at least $2^r$
nodes in its tree, and the height of any tree is $\Oh(\log n)$. Ranks of nodes along any
root path strictly increase.
\end{lemma}

\begin{proof}
By induction on the number of unions. A rank-$0$ root has its $1\ge 2^0$ node. A root reaches
rank $r$ only by uniting two rank-$(r-1)$ roots, each with $\ge 2^{r-1}$ nodes by induction,
so the merged tree has $\ge 2^r$ nodes. Hence $2^r\le n$, i.e. $r\le\log_2 n$. Path
compression only moves nodes closer to the root and never raises a rank, so it preserves the
node-count-per-rank bound. Ranks strictly increase from child to parent because a node is
attached only under a strictly-or-equally-ranked root, and the equal case bumps the parent's
rank. \end{proof}

So \emph{without} path compression each operation is already $\Oh(\log n)$. Path compression
pushes this down to $\Oh(\log^* n)$ (iterated logarithm) by a simpler argument, and to the
$\alpha(n)$ of Theorem~\ref{ch42:thm-uf} by Tarjan's full potential argument. The companion
fact to carry away:

\begin{complexitybox}
\textbf{Union-find.} Space $\Oh(n)$. With union by rank \emph{and} path compression,
$m$ operations in $\Oh(m\,\alpha(n))$ -- effectively $\Oh(1)$ each. With only one of the two
tricks, $\Oh(\log n)$ per operation. Naive (neither), $\Oh(n)$ worst case.
\end{complexitybox}

\begin{applicationbox}
Kruskal's MST sorts edges then does $m$ union-finds: the sort dominates at $\Oh(m\log m)$,
the connectivity bookkeeping is essentially free. Union-find also drives offline LCA
(Tarjan's algorithm), incremental connected components, cycle detection in undirected graphs,
and the merge step in Bor\r{u}vka's algorithm. Whenever a problem is ``maintain a partition
under merges,'' reach for it.
\end{applicationbox}

\section{Heaps and Priority Queues}
\label{ch42:heaps}

A priority queue stores keyed elements and supports \textsc{Insert}, \textsc{ExtractMin}
(remove and return the smallest key), and often \textsc{DecreaseKey}. Dijkstra and Prim
\emph{are} priority queues with a relaxation loop wrapped around them, so the heap's cost
profile directly sets the shortest-path and MST bounds.

\begin{intuitionbox}
A \emph{binary heap} is a complete binary tree stored in an array, with the heap order:
every node's key is $\le$ its children's. The minimum sits at the root, index $1$. Children
of $i$ live at $2i$ and $2i+1$; the parent of $i$ is $\lfloor i/2\rfloor$. To insert, drop
the new key at the end and \emph{sift up} swapping with the parent while it is smaller. To
extract the min, move the last element to the root and \emph{sift down} swapping with the
smaller child while it violates order. Both touch one root-to-leaf path: $\Oh(\log n)$.
\end{intuitionbox}

\begin{center}
\begin{tikzpicture}[
  every node/.style={circle,draw=nodeblue,fill=blue!6,minimum size=7mm,inner sep=0pt,font=\small},
  level distance=10.5mm,
  level 1/.style={sibling distance=26mm},
  level 2/.style={sibling distance=13mm},
  >=Stealth]
\node[fill=targetgreen!20,draw=targetgreen!80!black,font=\small\bfseries] {$2$}
  child {node {$5$} child {node {$9$}} child {node {$7$}}}
  child {node {$4$} child {node {$8$}} child {node {$6$}}};
\begin{scope}[xshift=6.0cm,yshift=-1.2cm,every node/.style={rectangle,draw=nodeblue,fill=blue!4,minimum size=7mm,inner sep=1pt,font=\small\ttfamily}]
\node (i1) at (0,0) {2};\node (i2) at (0.8,0) {5};\node (i3) at (1.6,0) {4};
\node (i4) at (2.4,0) {9};\node (i5) at (3.2,0) {7};\node (i6) at (4.0,0) {8};\node (i7) at (4.8,0) {6};
\foreach \n in {1,2,3,4,5,6,7} {\node[draw=none,fill=none,font=\tiny] at (0.8*\n-0.8,0.6) {\n};}
\node[draw=none,fill=none,font=\small] at (2.4,-0.8) {array layout (1-indexed)};
\end{scope}
\end{tikzpicture}
\end{center}

\begin{algorithm}[H]
\caption{Binary-heap sift operations (min-heap, array $A$, size $s$)}
\DontPrintSemicolon
\SetKwFunction{FUp}{SiftUp}\SetKwFunction{FDown}{SiftDown}\SetKwFunction{FIns}{Insert}\SetKwFunction{FEx}{ExtractMin}
\SetKwProg{Fn}{function}{:}{}
\Fn{\FUp{$i$}}{
  \While{$i>1$ \textbf{and} $A[\lfloor i/2\rfloor] > A[i]$}{
    swap $A[i]\leftrightarrow A[\lfloor i/2\rfloor]$;\quad $i\gets\lfloor i/2\rfloor$\;}
}
\Fn{\FDown{$i$}}{
  \While{$2i\le s$}{
    $j\gets 2i$;\quad \lIf{$j+1\le s$ \textbf{and} $A[j+1]<A[j]$}{$j\gets j+1$}
    \lIf{$A[i]\le A[j]$}{\textbf{break}}
    swap $A[i]\leftrightarrow A[j]$;\quad $i\gets j$\;}
}
\Fn{\FIns{$x$}}{$s\gets s+1$;\ $A[s]\gets x$;\ \FUp{$s$}\;}
\Fn{\FEx{}}{$m\gets A[1]$;\ $A[1]\gets A[s]$;\ $s\gets s-1$;\ \FDown{$1$};\ \Return $m$\;}
\end{algorithm}

\begin{whybox}
Correctness rests on a single invariant: \emph{after every operation, $A[\lfloor
i/2\rfloor]\le A[i]$ for all $i\ge 2$}. \textsc{SiftUp} restores it by repairing the one
edge that a fresh leaf could violate, walking the violation upward; \textsc{SiftDown} repairs
the one edge a displaced root could violate, walking it downward. Because the tree is
complete, its height is $\lfloor\log_2 n\rfloor$, so each sift does $\Oh(\log n)$ comparisons.
The root being the global minimum follows by transitivity along any root-to-node path.
\end{whybox}

\begin{labbox}
Build a heap by inserting $5,3,8,1,4$ into an empty min-heap.
\begin{itemize}
\item Insert $5$: $[5]$.
\item Insert $3$: $[5,3]\to$ sift up $\to[3,5]$.
\item Insert $8$: $[3,5,8]$, no swap.
\item Insert $1$: $[3,5,8,1]\to$ sift up: $1<5$ swap $\to[3,1,8,5]$; $1<3$ swap $\to[1,3,8,5]$.
\item Insert $4$: $[1,3,8,5,4]\to$ sift up: parent of index 5 is index 2 ($=3$), $4>3$, stop.
\end{itemize}
Now \textsc{ExtractMin}: returns $1$, moves $4$ to root $\to[4,3,8,5]$, sift down: smaller
child is $3$, swap $\to[3,4,8,5]$; children of index 2 are $5$ (index 4) only, $4<5$ stop.
Result $[3,4,8,5]$, min is $3$. The trace matches the invariant at every step.
\end{labbox}

\paragraph{Heapify in linear time.} Building a heap from an unsorted array of $n$ keys looks
like $n$ inserts at $\Oh(n\log n)$, but bottom-up \textsc{SiftDown} from index $\lfloor
n/2\rfloor$ down to $1$ is only $\Oh(n)$. The reason: a node at height $h$ costs $\Oh(h)$, and
there are $\le n/2^{h+1}$ nodes at height $h$, so the total is
\[
\sum_{h\ge 0}\frac{n}{2^{h+1}}\,\Oh(h)=\Oh\!\Big(n\sum_{h\ge 0}\frac{h}{2^h}\Big)=\Oh(2n)=\Oh(n),
\]
since $\sum_{h\ge0} h/2^h = 2$. Most nodes are leaves, where sifting down does nothing.

\begin{complexitybox}
\textbf{Heaps.} Binary heap: \textsc{Insert}, \textsc{ExtractMin}, \textsc{DecreaseKey} all
$\Oh(\log n)$; build $\Oh(n)$; space $\Oh(n)$. Binomial and pairing heaps trade constants.
\emph{Fibonacci heaps} achieve $\Oh(1)$ amortized \textsc{Insert} and \textsc{DecreaseKey}
with $\Oh(\log n)$ amortized \textsc{ExtractMin}, which is what gives Dijkstra its
$\Oh(m+n\log n)$ bound. In practice a binary or $d$-ary heap usually wins on constants.
\end{complexitybox}

\begin{connectionbox}
Williams introduced the binary heap in 1964 as the core structure of heapsort; Floyd gave the
linear-time bottom-up build the same year. Fibonacci heaps (Fredman and Tarjan, 1984) were
designed precisely to speed up Dijkstra's algorithm and minimum spanning trees, showing how
amortized potential functions absorb cascading cuts.
\end{connectionbox}

\section{Fenwick Trees (Binary Indexed Trees)}
\label{ch42:fenwick}

Maintain an array $a[1\dots n]$ under two operations: \emph{point update} $a[i]\mathrel{+}=v$
and \emph{prefix sum} $\sum_{j\le i} a[j]$. A plain array does updates in $\Oh(1)$ but
prefix sums in $\Oh(n)$; a prefix-sum array flips that. The Fenwick tree does both in
$\Oh(\log n)$ with a single array of size $n$ and a bit trick: isolating the lowest set bit.

\begin{intuitionbox}
Let $\lowbit(i)=i\,\&\,(-i)$ be the lowest set bit of $i$ (in two's complement, $-i$ flips
and adds one, so $i\,\&\,(-i)$ isolates the rightmost $1$). Store at index $i$ the partial
sum of the $\lowbit(i)$ entries ending at $i$, namely $a[i-\lowbit(i)+1\dots i]$. Then a
prefix sum up to $i$ is assembled by repeatedly subtracting $\lowbit$: each step jumps over a
block of size a power of two, and the blocks tile $[1,i]$ exactly. The number of jumps is the
number of set bits in $i$, at most $\log_2 n$.
\end{intuitionbox}

\begin{center}
\begin{tikzpicture}[>=Stealth,every node/.style={font=\small}]
\foreach \i in {1,...,8} {\draw (\i-0.5,0) rectangle (\i+0.5,0.6);\node at (\i,0.3){\i};}
\draw[targetgreen,thick] (0.5,-0.2)--(1.5,-0.2) node[right,black]{\scriptsize $T[1]{=}a[1]$};
\draw[sourceblue,thick] (0.5,-0.55)--(2.5,-0.55) node[right,black]{\scriptsize $T[2]{=}a[1..2]$};
\draw[targetgreen,thick] (2.5,-0.9)--(3.5,-0.9) node[right,black]{\scriptsize $T[3]{=}a[3]$};
\draw[quotientorange,thick] (0.5,-1.25)--(4.5,-1.25) node[right,black]{\scriptsize $T[4]{=}a[1..4]$};
\draw[targetgreen,thick] (4.5,-1.6)--(5.5,-1.6) node[right,black]{\scriptsize $T[5]{=}a[5]$};
\draw[sourceblue,thick] (4.5,-1.95)--(6.5,-1.95) node[right,black]{\scriptsize $T[6]{=}a[5..6]$};
\draw[targetgreen,thick] (6.5,-2.3)--(7.5,-2.3) node[right,black]{\scriptsize $T[7]{=}a[7]$};
\draw[structurepurple,thick] (0.5,-2.65)--(8.5,-2.65) node[right,black]{\scriptsize $T[8]{=}a[1..8]$};
\end{tikzpicture}
\end{center}

\begin{algorithm}[H]
\caption{Fenwick tree (array $T[1\dots n]$ initialized to $0$)}
\DontPrintSemicolon
\SetKwFunction{FUpd}{Update}\SetKwFunction{FQ}{PrefixSum}
\SetKwProg{Fn}{function}{:}{}
\Fn{\FUpd{$i,v$}}{
  \While{$i\le n$}{$T[i]\gets T[i]+v$;\quad $i\gets i+\lowbit(i)$\tcp*{climb to covering blocks}}
}
\Fn{\FQ{$i$}}{
  $s\gets 0$\;
  \While{$i>0$}{$s\gets s+T[i]$;\quad $i\gets i-\lowbit(i)$\tcp*{descend over tiling blocks}}
  \Return $s$\;
}
\end{algorithm}

A range sum $a[l\dots r]$ is $\FQ{r}-\FQ{l-1}$. By swapping the roles -- range update and
point query -- via a difference array, or running two Fenwick trees, you get range-update
range-sum, still $\Oh(\log n)$.

\begin{whybox}
Why do the descending jumps tile $[1,i]$ exactly? Index $i$ covers
$(i-\lowbit(i),\,i]$. Set $i'=i-\lowbit(i)$; this clears the lowest set bit, so $i'$ covers a
block ending exactly where the previous block began, with no gap and no overlap. Iterating
clears set bits one at a time until $i$ reaches $0$, so the visited blocks partition $[1,i]$.
Symmetrically, $i+\lowbit(i)$ in \textsc{Update} walks to every block whose range contains
$i$, which are exactly the cells that must change. Both loops do one iteration per set bit
encountered, bounded by $\lfloor\log_2 n\rfloor+1$.
\end{whybox}

\begin{labbox}
Let $n=8$, all zero. \FUpd{$3,5$}: visits $3\,(\lowbit 1)\to4\,(\lowbit 4)\to8\,(\lowbit
8)\to16>8$ stop. So $T[3],T[4],T[8]\mathrel{+}=5$. Now \FQ{$6$}: visits
$6\,(+T[6]{=}0)\to4\,(+T[4]{=}5)\to0$ stop, returns $5$. Correct: only $a[3]=5$ lies in
$[1,6]$. \FQ{$2$} visits $2\to0$, returns $0$. Range sum $a[3..6]=\FQ{6}-\FQ{2}=5-0=5$. The
bit jumps did three array touches instead of scanning six cells.
\end{labbox}

\section{Segment Trees}
\label{ch42:segtree}

The Fenwick tree is a specialist: fast, tiny, but tied to invertible aggregates like sums.
The \emph{segment tree} is the generalist. It answers ``aggregate over $a[l\dots r]$'' for
\emph{any} associative operation -- min, max, gcd, sum, matrix product -- and supports point
and (with lazy propagation) range updates, all in $\Oh(\log n)$.

\begin{intuitionbox}
Build a balanced binary tree over the index range $[1,n]$. The root owns $[1,n]$; each
internal node owning $[l,r]$ splits at $m=\lfloor(l+r)/2\rfloor$ into children $[l,m]$ and
$[m+1,r]$; leaves own single indices. Each node caches the aggregate of its range. A query
$[ql,qr]$ recurses, returning a node's cached value whole when its range lies inside
$[ql,qr]$, and recursing only when the query straddles the split. At most $\Oh(\log n)$ nodes
are ever fully inside while their parent straddles, so the query visits $\Oh(\log n)$ nodes.
\end{intuitionbox}

\begin{figure}[H]
\centering
\begin{tikzpicture}[
  every node/.style={draw,rounded corners=2pt,minimum width=11mm,minimum height=5mm,inner sep=2pt,font=\scriptsize},
  canonical/.style={draw=trapred!90!black,fill=quotientorange!30,line width=1.3pt},
  level distance=10mm,>=Stealth,
  level 1/.style={sibling distance=52mm},
  level 2/.style={sibling distance=26mm},
  level 3/.style={sibling distance=13mm}
]
\node[fill=gray!10] {$[1,8]{:}\,29$}
 child {node[fill=gray!10] {$[1,4]{:}\,12$}
   child {node[fill=gray!10] {$[1,2]{:}\,7$} 
     child {node[fill=gray!10] {$[1]{:}3$}} 
     child {node[canonical] {$[2]{:}4$}}}
   child {node[canonical] {$[3,4]{:}\,5$} 
     child {node[fill=gray!10] {$[3]{:}1$}} 
     child {node[fill=gray!10] {$[4]{:}4$}}}}
 child {node[fill=gray!10] {$[5,8]{:}\,17$}
   child {node[canonical] {$[5,6]{:}\,9$} 
     child {node[fill=gray!10] {$[5]{:}6$}} 
     child {node[fill=gray!10] {$[6]{:}3$}}}
   child {node[fill=gray!10] {$[7,8]{:}\,8$} 
     child {node[canonical] {$[7]{:}5$}} 
     child {node[fill=gray!10] {$[8]{:}3$}}}};
\end{tikzpicture}
\caption{Segment tree range query decomposition for query range $[2, 7]$ on array $a=[3,4,1,4,6,3,5,3]$. The query recursively traverses the tree and collects the maximal disjoint canonical subtrees completely contained in $[2, 7]$ (highlighted in orange with red borders): nodes $[2]{=}4$, $[3,4]{=}5$, $[5,6]{=}9$, and $[7]{=}5$. Summing these $4$ cached values immediately yields $\sum a[2\dots 7] = 4+5+9+5 = 23$ in $\Oh(\log n)$ time without touching all $6$ individual leaf elements.}
\label{ch42:fig-segtree-query}
\end{figure}
The leaves hold $a=[3,4,1,4,6,3,5,3]$; each internal node holds the sum of its subtree. A
query for $\sum a[2\dots7]$ stitches together $[2]{=}4$, $[3,4]{=}5$, $[5,6]{=}9$, $[7]{=}5$,
totaling $23$ -- four cached nodes, no full scan.

\begin{algorithm}[H]
\caption{Segment tree: build and range query (sum; node $p$ owns $[l,r]$)}
\DontPrintSemicolon
\SetKwFunction{FBuild}{Build}\SetKwFunction{FQ}{Query}\SetKwFunction{FUpd}{Update}
\SetKwProg{Fn}{function}{:}{}
\Fn{\FBuild{$p,l,r$}}{
  \lIf{$l=r$}{$t[p]\gets a[l]$; \Return}
  $m\gets\lfloor(l+r)/2\rfloor$\;
  \FBuild{$2p,l,m$};\quad \FBuild{$2p{+}1,m{+}1,r$}\;
  $t[p]\gets t[2p]+t[2p{+}1]$\;
}
\Fn{\FQ{$p,l,r,ql,qr$}}{
  \lIf{$qr<l$ \textbf{or} $r<ql$}{\Return $0$ \tcp*[f]{disjoint: identity}}
  \lIf{$ql\le l$ \textbf{and} $r\le qr$}{\Return $t[p]$ \tcp*[f]{fully inside}}
  $m\gets\lfloor(l+r)/2\rfloor$\;
  \Return \FQ{$2p,l,m,ql,qr$} $+$ \FQ{$2p{+}1,m{+}1,r,ql,qr$}\;
}
\Fn{\FUpd{$p,l,r,i,v$}}{
  \lIf{$l=r$}{$t[p]\gets v$; \Return}
  $m\gets\lfloor(l+r)/2\rfloor$\;
  \lIf{$i\le m$}{\FUpd{$2p,l,m,i,v$}}\lElse{\FUpd{$2p{+}1,m{+}1,r,i,v$}}
  $t[p]\gets t[2p]+t[2p{+}1]$\;
}
\end{algorithm}

\begin{theorem}[Query cost]
\label{ch42:thm-segtree}
A range query visits $\Oh(\log n)$ nodes.
\end{theorem}

\begin{proof}
At each depth of the tree, consider the nodes the recursion enters. A node is entered only if
its range overlaps $[ql,qr]$ partially (a fully-inside or fully-outside node returns without
recursing). Claim: at most two nodes per depth are entered with a \emph{partial} overlap --
one containing the left endpoint $ql$, one containing $qr$. Any node strictly between them is
fully inside and stops; any node outside is disjoint and stops. The two boundary nodes each
spawn at most two children at the next depth, but again only the boundary-touching ones
recurse. So the number of recursing nodes per level is $\Oh(1)$, and the tree has depth
$\Oh(\log n)$, giving $\Oh(\log n)$ visited nodes total. \end{proof}

\begin{trapbox}
The classic bug is an off-by-one in the disjoint test. Writing $qr<l$ \emph{or} $r<ql$ is
correct; writing $qr\le l$ accidentally drops the boundary index when $ql=qr=l$. Equally
common: forgetting that the identity element depends on the operation -- $0$ for sum,
$+\infty$ for min, $-\infty$ for max, the empty product for products. Return the wrong
identity on the disjoint branch and a single query silently corrupts the whole answer.
\end{trapbox}

\subsection{Lazy propagation}

To support \emph{range} updates (add $v$ to all of $a[l\dots r]$) without touching every leaf,
store a pending \emph{lazy} value at each node: ``this whole subtree still owes this update.''
Apply it to the node's cached aggregate immediately, but push it down to children only when
the recursion actually needs to descend past this node. A node fully inside the update range
absorbs the change in $\Oh(1)$ and parks the lazy tag; the $\Oh(\log n)$ boundary argument of
Theorem~\ref{ch42:thm-segtree} still bounds the work.

\begin{whybox}
The lazy tag is correct because the cached aggregate is kept consistent with all applied-but-
not-pushed updates: a range-add of $v$ over a node owning $k$ indices raises its sum by
$k\cdot v$ instantly, and the children's caches stay valid relative to \emph{their} own tags.
Pushing down before recursing guarantees a child is never read while an ancestor's tag is
still owed to it. This is the same ``pay lazily, settle on demand'' idea as the
amortized accounting method, applied to subtrees instead of to a sequence of operations.
\end{whybox}

\begin{complexitybox}
\textbf{Range structures.} Fenwick: $\Oh(\log n)$ per op, $\Oh(n)$ space, invertible
aggregates only, smallest constants. Segment tree: $\Oh(\log n)$ per op, $\Oh(n)$ space
(array of size $\le 4n$), any associative aggregate, range updates via lazy propagation.
Sparse table (static, idempotent ops like min): $\Oh(n\log n)$ build, $\Oh(1)$ query, no
updates.
\end{complexitybox}

\begin{connectionbox}
The tree chapters rely on these structures throughout. Heavy-light decomposition
(Chapter~19) and the Euler-tour technique (Chapter~21) reduce path and subtree queries to
\emph{range} queries on an array, then hand them to a segment tree. The LCA-as-RMQ reduction
(Chapter~23) feeds a sparse table or segment tree. Link-cut trees (Chapter~22) generalize the
segment-tree idea to a forest that changes shape. If a tree algorithm in this book runs in
$\Oh(n\log n)$, a range structure is almost certainly why.
\end{connectionbox}

\section{Linear Algebra for the Spectral Chapters}
\label{ch42:linalg}

The Laplacian, spectral-clustering, expander, and electrical-flow chapters all speak the
language of symmetric matrices and their eigenvalues. We collect the facts they use. This is
a summary, not a course; the companion \emph{Linear Algebra} volume develops everything
below from scratch, and we cite it freely.

\subsection{Symmetric matrices and the spectral theorem}

A real matrix $\mathbf{M}\in\R^{n\times n}$ is \emph{symmetric} if $\mathbf{M}^\top=\mathbf
{M}$. Adjacency matrices of undirected graphs and graph Laplacians are symmetric, which is
exactly why graphs have real spectra. A nonzero vector $x$ is an \emph{eigenvector} with
\emph{eigenvalue} $\lambda$ if $\mathbf{M}x=\lambda x$: the matrix merely stretches $x$ by
$\lambda$ without rotating it.

\begin{theorem}[Spectral theorem, real symmetric case]
\label{ch42:thm-spectral}
Every symmetric $\mathbf{M}\in\R^{n\times n}$ has $n$ real eigenvalues
$\lambda_1\le\lambda_2\le\dots\le\lambda_n$ (with multiplicity) and an associated
\emph{orthonormal} basis of eigenvectors $u_1,\dots,u_n$. Equivalently
\[
\mathbf{M}=\sum_{i=1}^n \lambda_i\,u_i u_i^\top = \mathbf{U}\Lambda\mathbf{U}^\top,
\]
where $\mathbf{U}=[u_1\,\cdots\,u_n]$ is orthogonal ($\mathbf{U}^\top\mathbf{U}=\mathbf{I}$)
and $\Lambda=\diag(\lambda_1,\dots,\lambda_n)$.
\end{theorem}

We do not prove it here -- see the companion \emph{Linear Algebra} volume -- but the consequences run through
every spectral chapter. The eigenvectors give an orthonormal coordinate system in which
$\mathbf{M}$ acts diagonally; problems that look coupled in the standard basis decouple into
$n$ independent one-dimensional problems, one per eigenvalue.

\begin{intuitionbox}
Picture $\mathbf{M}$ as a deformation of space. The spectral theorem says a symmetric
deformation has $n$ perpendicular ``principal axes'' (the $u_i$) along which it is pure
stretch by $\lambda_i$, no shear. The Laplacian's smallest eigenvalue is always $0$ with
eigenvector the all-ones vector, because $\Lap\mathbf{1}=0$; the \emph{number} of zero
eigenvalues counts connected components, and the second-smallest $\lambda_2$ -- the Fiedler
value -- measures how hard the graph is to cut. That single number is central to Chapters~33
and~34 (spectral clustering and expanders).
\end{intuitionbox}

\begin{figure}[H]
\centering
\begin{tikzpicture}[>=Stealth,scale=0.9]
  % Panel A: Spectral Mapping of Unit Circle to Ellipse
  \begin{scope}[shift={(0,0)}]
    \node[font=\footnotesize\scshape,text=sourceblue!90!black] at (2.1,2.4) {Spectral Mapping $\mathbf{M} = \mathbf{U}\mathbf{\Lambda}\mathbf{U}^\top$};
    % Unit circle
    \draw[gray!40,thick] (0,0) circle (1.15);
    \draw[->,sourceblue,very thick] (0,0) -- (30:1.15) node[above right=1pt,font=\scriptsize]{$u_1$};
    \draw[->,targetgreen,very thick] (0,0) -- (120:1.15) node[above left=1pt,font=\scriptsize]{$u_2$};
    \node[font=\tiny,text=gray!75!black] at (0,-1.95) {Unit circle $\|x\|=1$};

    % Transformation arrow
    \draw[->,very thick,gray!60] (1.6,0) -- node[above=2pt,font=\tiny\bfseries]{$\mathbf{M}$} (2.6,0);

    % Transformed Ellipse
    \begin{scope}[shift={(4.2,0)}]
      \draw[gray!40,thick,rotate=30] (0,0) ellipse (0.68 and 1.55);
      \draw[->,sourceblue,very thick] (0,0) -- (30:0.68) node[below right=1pt,font=\scriptsize]{$\lambda_1 u_1$};
      \draw[->,targetgreen,very thick] (0,0) -- (120:1.55) node[above left=1pt,font=\scriptsize]{$\lambda_2 u_2$};
      \node[font=\tiny,text=gray!75!black] at (0,-1.95) {Principal axes stretched by $\lambda_i$};
    \end{scope}
  \end{scope}

  % Panel B: Quadratic Form Level Sets & Energy Bowl
  \begin{scope}[shift={(9.0,0)}]
    \node[font=\footnotesize\scshape,text=sourceblue!90!black] at (0.3,2.4) {PSD Energy Bowl $q(x) = x^\top \mathbf{M} x$};
    % Elliptical contour level sets (wider along flatter u1 direction 30 deg, narrower along steeper u2 direction 120 deg)
    \draw[sourceblue!30,thick,rotate=30] (0,0) ellipse (0.65 and 0.38);
    \draw[sourceblue!55,thick,rotate=30] (0,0) ellipse (1.15 and 0.68);
    \draw[sourceblue!80,thick,rotate=30] (0,0) ellipse (1.65 and 0.98);
    
    % Eigenvector directions
    \draw[dashed,gray!60] (-1.85,-1.07) -- (1.85,1.07) node[right=2pt,font=\tiny]{$u_1$ ($\lambda_1$, flatter)};
    \draw[dashed,gray!60] (0.7,-1.21) -- (-0.7,1.21) node[above left=1pt,font=\tiny]{$u_2$ ($\lambda_2$, steeper)};
    
    \node[fill=trapred,circle,inner sep=1.5pt] (orig) at (0,0) {};
    \draw[<-,trapred!85!black,thin] (0.08,-0.08) -- (0.85,-1.35)
      node[right=1pt,font=\tiny\bfseries,text=trapred]{minimum $q(0)=0$};

    \node[font=\tiny,text=gray!75!black] at (0.3,-1.95) {Level curves $x^\top \mathbf{M} x = c$;\quad $R_{\mathbf{M}}(x) \in [\lambda_1, \lambda_n]$};
  \end{scope}
\end{tikzpicture}
\caption{Geometry of symmetric and positive semidefinite matrices. Left: the Spectral Theorem states that a symmetric matrix $\mathbf{M}$ acts by pure stretch along mutually orthogonal eigenvectors $u_1, u_2$ without rotation or shearing, mapping the unit circle to an aligned ellipse. Right: when $\mathbf{M} \succeq 0$, the quadratic energy surface $q(x) = x^\top \mathbf{M} x \ge 0$ forms a convex parabolic bowl whose concentric level sets are ellipsoids, with curvature $\lambda_i$ along direction $u_i$. The Rayleigh quotient $R_{\mathbf{M}}(x) = \frac{x^\top \mathbf{M} x}{\|x\|^2}$ smoothly interpolates between minimum curvature $\lambda_1$ and maximum curvature $\lambda_n$.}
\label{ch42:fig-spectral-geometry}
\end{figure}
A symmetric $\mathbf{M}$ maps the unit circle to an ellipse whose axes are the orthonormal
eigenvectors $u_i$, each scaled by its eigenvalue $\lambda_i$. No rotation of the axes
occurs; that is what ``symmetric means real spectrum, orthogonal eigenvectors'' looks like
geometrically.

\subsection{Positive semidefiniteness and the Rayleigh quotient}

\begin{definition}[PSD and PD]
\label{ch42:def-psd}
A symmetric $\mathbf{M}$ is \emph{positive semidefinite} (PSD), written $\mathbf{M}\succeq 0$,
if $x^\top\mathbf{M}x\ge 0$ for all $x\in\R^n$; it is \emph{positive definite} (PD,
$\mathbf{M}\succ 0$) if the inequality is strict for all $x\ne 0$.
\end{definition}

\begin{proposition}[Equivalent characterizations of PSD]
\label{ch42:prop-psd}
For symmetric $\mathbf{M}$ the following are equivalent: (i) $\mathbf{M}\succeq 0$;
(ii) every eigenvalue $\lambda_i\ge 0$; (iii) $\mathbf{M}=\mathbf{B}^\top\mathbf{B}$ for some
matrix $\mathbf{B}$.
\end{proposition}

\begin{proof}[Sketch]
(i)$\Rightarrow$(ii): take $x=u_i$, then $u_i^\top\mathbf{M}u_i=\lambda_i\ge 0$.
(ii)$\Rightarrow$(iii): with $\lambda_i\ge0$ set $\mathbf{B}=\Lambda^{1/2}\mathbf{U}^\top$, so
$\mathbf{B}^\top\mathbf{B}=\mathbf{U}\Lambda\mathbf{U}^\top=\mathbf{M}$.
(iii)$\Rightarrow$(i): $x^\top\mathbf{B}^\top\mathbf{B}x=\norm{\mathbf{B}x}^2\ge 0$.
\end{proof}

The graph Laplacian is the headline example: $\Lap=\sum_{(u,v)\in E} (e_u-e_v)(e_u-e_v)^\top$,
which is manifestly of the form $\mathbf{B}^\top\mathbf{B}$ with $\mathbf{B}$ the signed
incidence matrix, hence PSD. Concretely, $x^\top\Lap x=\sum_{(u,v)\in E}(x_u-x_v)^2\ge 0$,
the ``Dirichlet energy'' minimized in Chapters~31, 33, 35, and~41.

\begin{definition}[Rayleigh quotient]
\label{ch42:def-rayleigh}
For symmetric $\mathbf{M}$ and $x\ne 0$, the \emph{Rayleigh quotient} is
\[
R_{\mathbf{M}}(x)=\frac{x^\top\mathbf{M}x}{x^\top x}.
\]
\end{definition}

\begin{theorem}[Variational characterization of eigenvalues]
\label{ch42:thm-courant}
With $\lambda_1\le\dots\le\lambda_n$,
\[
\lambda_1=\min_{x\ne 0} R_{\mathbf{M}}(x),\qquad
\lambda_n=\max_{x\ne 0} R_{\mathbf{M}}(x),
\]
and more generally (Courant--Fischer)
\[
\lambda_k=\min_{\dim S=k}\ \max_{0\ne x\in S} R_{\mathbf{M}}(x).
\]
\end{theorem}

\begin{proof}
Expand $x=\sum_i c_i u_i$ in the orthonormal eigenbasis. Then $x^\top x=\sum_i c_i^2$ and
$x^\top\mathbf{M}x=\sum_i \lambda_i c_i^2$, so
\[
R_{\mathbf{M}}(x)=\frac{\sum_i \lambda_i c_i^2}{\sum_i c_i^2}
\]
is a weighted average of the eigenvalues with nonnegative weights $c_i^2$. A weighted average
lies between the smallest and largest values, attaining $\lambda_1$ at $x=u_1$ and $\lambda_n$
at $x=u_n$.

For Courant--Fischer, prove both inequalities. Taking $S_0=\mathrm{span}\{u_1,\dots,u_k\}$,
every $x\in S_0$ averages only $\lambda_1,\dots,\lambda_k$, so $\max_{x\in S_0}R_{\mathbf{M}}(x)=\lambda_k$
and the minimum over subspaces is at most $\lambda_k$. Conversely, let $S$ be any subspace with
$\dim S=k$, and let $U=\mathrm{span}\{u_k,\dots,u_n\}$, of dimension $n-k+1$. The dimensions add to
$n+1>n$, so $S\cap U$ contains a nonzero $x$. That $x$ averages only $\lambda_k,\dots,\lambda_n$, so
$R_{\mathbf{M}}(x)\ge\lambda_k$, and the maximum over $S$ is at least $\lambda_k$. \end{proof}

\begin{whybox}
Much of spectral graph theory starts from this theorem. It turns ``find the second eigenvalue
$\lambda_2$ of the Laplacian'' into ``minimize $\sum_{(u,v)\in E}(x_u-x_v)^2$ over unit
vectors $x$ orthogonal to $\mathbf{1}$.'' That optimization \emph{is} the relaxed balanced-cut
problem: assign each vertex a real number so that adjacent vertices land close, subject to a
spread constraint. After degree normalization (using the normalized Laplacian, whose
eigenvalues lie in $[0,2]$), Cheeger's inequality sandwiches the conductance between
$\lambda_2/2$ and $\sqrt{2\lambda_2}$, which is how the Fiedler eigenvector becomes a graph
partitioner. Every appearance of ``the energy of a vector'' or ``the Rayleigh quotient'' in
the later chapters traces back to Theorem~\ref{ch42:thm-courant}.
\end{whybox}

\begin{example}[Trace and determinant from the spectrum]
\label{ch42:ex-trace}
Two identities used repeatedly: $\tr(\mathbf{M})=\sum_i \lambda_i$ and
$\det(\mathbf{M})=\prod_i\lambda_i$. The first counts, e.g., closed walks: $\tr(\Adjm^k)$
equals both the number of closed walks of length $k$ and $\sum_i \lambda_i^k$. The second
underlies the matrix-tree theorem: the number of spanning trees equals
$\frac{1}{n}\prod_{i\ge 2}\lambda_i$ over the nonzero Laplacian eigenvalues.
\end{example}

\begin{connectionbox}
Pointers into the main text. Chapter~31 builds $\Lap$ and proves the component count equals
$\dim\ker\Lap$. Chapter~32 develops the full spectrum and the matrix-tree theorem.
Chapter~33 runs $k$-means on the bottom Fiedler vectors. Chapter~34 ties $\lambda_2$ to
expansion. Chapter~35 reads off effective resistance from the pseudoinverse $\Lap^+$, and
Chapter~41's Laplacian solvers approximate $\Lap^{+}b$ in near-linear time. For derivations
of the spectral theorem, PSD factorizations, pseudoinverses, and Courant--Fischer, see the
companion \emph{Linear Algebra} volume, in the chapters on symmetric operators and quadratic
forms.
\end{connectionbox}

\section{Probability for the Randomized Chapters}
\label{ch42:probability}

Karger's contraction algorithm, randomized min-cut, random walks, PageRank, and the
streaming sketches all rely on a small amount of discrete probability, collected here.

\subsection{Expectation, linearity, and variance}

A discrete random variable $X$ takes values $x$ with probabilities $\Pr[X=x]$. Its
\emph{expectation} is $\E[X]=\sum_x x\,\Pr[X=x]$. The single most-used fact is
\emph{linearity of expectation}: for \emph{any} random variables (independent or not) and
constants,
\[
\E\Big[\sum_i a_i X_i\Big]=\sum_i a_i\,\E[X_i].
\]
No independence needed -- this is what makes it so useful. Counting arguments that look
hopeless (``expected number of edges crossing a random cut'') collapse to summing one
indicator at a time.

\begin{example}[Indicator trick]
\label{ch42:ex-indicator}
To count the expected number of events that occur, write the count as a sum of indicator
variables $X_i=\mathbf{1}[\text{event }i]$, note $\E[X_i]=\Pr[\text{event }i]$, and sum. For a
uniformly random cut of a graph with $m$ edges, each edge is cut with probability $1/2$, so
the expected number of cut edges is $m/2$ -- the one-line proof that a max cut of size
$\ge m/2$ always exists.
\end{example}

The \emph{variance} $\Var[X]=\E[(X-\E[X])^2]=\E[X^2]-\E[X]^2$ measures spread, and for
\emph{independent} $X,Y$, $\Var[X+Y]=\Var[X]+\Var[Y]$.

\subsection{Concentration: from average to almost-sure}

Expectation tells you the center; concentration inequalities tell you a random variable
rarely strays far. Three, in increasing strength and assumption:

\begin{theorem}[Markov, Chebyshev, Chernoff]
\label{ch42:thm-concentration}
Let $X\ge 0$. \emph{(Markov)} $\Pr[X\ge t]\le \E[X]/t$. If $X$ has finite variance,
\emph{(Chebyshev)} $\Pr[\,\abs{X-\E[X]}\ge t\,]\le \Var[X]/t^2$. If $X=\sum_{i=1}^n X_i$ is a
sum of independent $\{0,1\}$ variables with $\mu=\E[X]$, then \emph{(Chernoff)} for
$0<\delta\le 1$,
\[
\Pr[\,\abs{X-\mu}\ge \delta\mu\,]\le 2\exp\!\Big(-\frac{\delta^2\mu}{3}\Big).
\]
\end{theorem}

\begin{proof}[Markov, and the bounds from it]
Since $X\ge 0$, $\E[X]=\sum_x x\Pr[X=x]\ge\sum_{x\ge t} x\Pr[X=x]\ge t\,\Pr[X\ge t]$,
giving Markov. Chebyshev is Markov applied to $(X-\E[X])^2$ with threshold $t^2$. Chernoff
applies Markov to the moment generating function $e^{sX}$ and optimizes over $s$; the
exponential decay is what makes ``repeat and take the majority'' so effective. \end{proof}

\begin{intuitionbox}
The qualitative ladder: Markov needs only nonnegativity and decays like $1/t$ -- weak but
universal. Chebyshev needs variance and decays like $1/t^2$ -- the basis of the median trick.
Chernoff needs independence and decays \emph{exponentially} -- this is why a randomized
algorithm with success probability $2/3$ can be boosted to $1-2^{-k}$ by running it $\Oh(k)$
times and voting. The error shrinks geometrically in the number of repetitions.
\end{intuitionbox}

\begin{center}
\begin{tikzpicture}
\begin{axis}[width=10cm,height=5.2cm,domain=1:6,samples=80,
  xlabel={deviation $t$ (in units of the scale)},ylabel={tail bound},
  ymin=0,ymax=1.05,legend pos=north east,legend style={font=\scriptsize},
  every axis plot/.append style={thick},tick label style={font=\scriptsize},
  label style={font=\scriptsize}]
\addplot[sourceblue]{1/x};            \addlegendentry{Markov $\sim 1/t$}
\addplot[targetgreen]{1/(x*x)};       \addlegendentry{Chebyshev $\sim 1/t^2$}
\addplot[trapred]{exp(-x*x/3)};       \addlegendentry{Chernoff $\sim e^{-t^2/3}$}
\end{axis}
\end{tikzpicture}
\end{center}
The three bounds in increasing strength: Markov decays polynomially and slowly, Chebyshev
faster, and Chernoff drops off exponentially. The price for the stronger bound is a stronger
hypothesis -- nonnegativity, then finite variance, then independence.

\begin{applicationbox}
Karger's algorithm contracts a random edge repeatedly; a fixed min cut survives all
contractions with probability $\ge 2/n^2$ (a linearity-and-conditioning computation). Running
it $\Oh(n^2\log n)$ times makes the failure probability $(1-2/n^2)^{\Th(n^2\log n)}\le 1/n$ by
the inequality $1-x\le e^{-x}$ -- pure Chernoff-style boosting. Chapters~35 and~36
(random walks and PageRank) use the same machinery to bound mixing times: a walk's distribution converges to
stationary at a rate governed by the spectral gap $1-\lambda_2$ of the transition matrix,
where here $\lambda_2$ is its second-largest eigenvalue,
tying the probability section straight back to the linear algebra one.
\end{applicationbox}

\begin{trapbox}
Independence is not optional in Chernoff. Linearity of expectation survives arbitrary
dependence, but variance addition and the exponential tail bound do \emph{not}: highly
correlated indicators can have the right mean and still pile all their mass at the extremes.
When the randomized chapters invoke Chernoff, check that the summands really are independent
(or use a martingale/Azuma variant when they are not). Mixing up ``$\E$ behaves nicely'' with
``the sum concentrates'' is the most common error in randomized analysis.
\end{trapbox}

\section{Historical Context, Controversies, and the Unified Curriculum}

\begin{historybox}
Each of the mathematical and algorithmic tools in this appendix grew out of a concrete problem in data organization, continuous physics, or statistical inference. Asymptotic notation began with Paul Bachmann ($1894$) and Edmund Landau in number theory, before Donald Knuth imported $\Oh$, $\Om$, and $\Th$ in the $1970$s to give algorithm analysis a standard classification system. The analysis of Disjoint-Set Union (DSU) produced one of the most famous surprises in theoretical computer science. Introduced by Bernard Galler and Michael Fischer in $1964$, the runtime of union-find with path compression remained open for a decade until Robert Tarjan published his $1975$ analysis. Tarjan proved that $m$ operations on $n$ elements require $\Theta(m\,\alpha(m,n))$ time, where $\alpha$ is the functional inverse of Ackermann's hyper-exponential growth function -- showing that the cost per operation is provably non-constant yet grows so slowly that $\alpha(n) \le 4$ for all $n$ within physical universe bounds. Priority queues saw a parallel evolution: J.~W.~J.~Williams introduced the binary heap in $1964$, Robert Floyd devised the linear-time bottom-up build that same year, and Michael Fredman and Robert Tarjan ($1984$) created Fibonacci heaps to accelerate Dijkstra's and Prim's algorithms via amortized potential functions.

Range query data structures emerged from geometric computing and data compression. In $1977$, Jon Bentley introduced segment trees to solve Klee's measure problem in computational geometry, pioneering the canonical decomposition of arbitrary intervals into logarithmic dyadic nodes. In $1994$, Peter Fenwick introduced the Binary Indexed Tree while developing arithmetic data compression algorithms, discovering that low-bit manipulation in two's-complement arithmetic creates a tree structure with zero pointer overhead. In parallel, the spectral foundations trace to Ernst Fischer ($1905$) and Richard Courant ($1920$), who formulated the Courant--Fischer min-max theorem to characterize eigenvalues as solutions to variational Rayleigh quotient optimizations over subspaces.

Several of these results corrected widespread misconceptions. Before Tarjan, algorithm designers were divided between believing path compression achieved $\Oh(1)$ amortized time or suspecting an unyielding $\Om(\log\log n)$ lower bound; both were wrong, and the true bound turned out to lie outside the elementary functions. When random matrices emerged in spectral graph theory, researchers routinely made the error of applying scalar Chernoff bounds element-by-element to matrices and taking a union bound over $n^2$ entries, which gives loose bounds, until Joel Tropp ($2012$) developed Matrix Chernoff bounds via non-commutative operator inequalities. What remains open includes developing truly cache-oblivious range query structures that match pointer-free theoretical bounds without empirical constant overhead, and simplifying non-commutative matrix concentration inequalities for graph applications.
\end{historybox}

\subsection{Summary of the Book and Companion Volumes}
\begin{connectionbox}
\textbf{What This Book Has Built}:
Across forty-two chapters, this book has covered graph algorithms from the basics up:
\begin{enumerate}
  \item \emph{Discrete Foundations}: From elementary traversals (BFS, DFS) and shortest paths (Dijkstra, Bellman--Ford) to trees, topological sorting, and disjoint sets.
  \item \emph{Combinatorial Optimization}: Maximum flows, minimum cuts, bipartite matching, link-cut trees, and heavy-light decomposition.
  \item \emph{NP-Hardness and Approximation}: Karp's web of reductions, Christofides' metric TSP, and Goemans--Williamson semidefinite hyperplane rounding.
  \item \emph{Dynamic and Streaming Algorithms}: Real-time forest maintenance via HdLT level hierarchies and sublinear-space processing via AGM linear sketches.
  \item \emph{Continuous and Spectral Methods}: The Laplacian matrix, Cheeger's isoperimetric inequality, spectral clustering, random walk electrical reciprocity, and the Spielman--Teng nearly-linear time Laplacian solver.
\end{enumerate}

\textbf{Where These Tools Go Next in the 8-Book Curriculum}:
\begin{itemize}
  \item In \textbf{Book 3 (Linear Algebra From the Ground Up)}, the eigenvalues, singular values, and Rayleigh quotients introduced here are developed from geometric first principles as linear transformations.
  \item In \textbf{Book 4 (Probability and Statistics From the Ground Up)}, the Markov chains, stationary distributions, and Chernoff concentration inequalities used in PageRank and random walks are explored in depth.
  \item In \textbf{Book 5 (Optimization From the Ground Up)}, the semidefinite programs, multiplicative weights updates, and interior-point methods that broke the max-flow barrier become our central study.
\end{itemize}
\end{connectionbox}

\section{Exercises}
\label{ch42:exercises}

\begin{exercise}[Asymptotic sorting]
Order these by growth rate, grouping any that are $\Th$ of each other:
$n\log n$, $2^{\log_2 n}$, $n^{1/\log_2 n}$, $\log(n!)$, $n^{1.001}$, $(\log n)^{\log n}$,
$2^{\sqrt{\log n}}$. (Hint: $\log(n!)=\Th(n\log n)$ by Stirling; $n^{1/\log_2 n}=2$.)
\end{exercise}

\begin{exercise}[Potential method]
A stack supports \textsc{Push}, \textsc{Pop}, and \textsc{MultiPop}$(k)$ which pops
$\min(k,\text{size})$ elements. Using the potential $\Phi=$ stack size, prove every operation
has $\Oh(1)$ amortized cost, so any sequence of $n$ operations on an initially empty stack
runs in $\Oh(n)$.
\end{exercise}

\begin{exercise}[Union-find lower bound by example]
Show that \emph{without} union by rank, path compression alone can produce a single $\find$
costing $\Om(n)$: give an adversarial sequence of unions building a path of length $n$, then
exhibit the expensive find. Then show union by rank alone (no compression) keeps every find
$\Oh(\log n)$, citing Lemma~\ref{ch42:lem-rank}.
\end{exercise}

\begin{exercise}[Heapify is tight]
Prove the bottom-up build is $\Th(n)$, not merely $\Oh(n)$: exhibit an input on which
$\Om(n)$ swaps actually occur. Contrast with the top-down ``insert $n$ times'' build and give
an input forcing it to do $\Th(n\log n)$ work.
\end{exercise}

\begin{exercise}[Fenwick from scratch]
Implement a Fenwick tree supporting range-add and range-sum using two Fenwick trees over a
difference array. Prove the formula $\sum_{j=1}^i a[j] = (i+1)\cdot B_1.\textsc{query}(i) -
B_2.\textsc{query}(i)$ for the auxiliary arrays $B_1,B_2$ you maintain, and verify it on
$n=4$ with one range update by hand.
\end{exercise}

\begin{exercise}[Segment tree, non-sum aggregate]
Adapt the segment tree of Section~\ref{ch42:segtree} to answer range-\emph{minimum} queries
with point updates. State the identity element, rewrite the combine step, and trace a query
for $\min a[2\dots7]$ on the array $[3,4,1,4,6,3,5,3]$, listing the cached nodes stitched
together.
\end{exercise}

\begin{exercise}[Break the lazy tag]
Here is a buggy lazy segment tree for range-add/range-sum: it pushes the lazy tag to children
\emph{after} recursing instead of before. Construct a sequence of two overlapping range-adds
and one query that returns the wrong sum, and explain which invariant the misordering
violates.
\end{exercise}

\begin{exercise}[Rayleigh quotient by hand]
Let $\mathbf{M}=\begin{psmallmatrix}2&1\\1&2\end{psmallmatrix}$. Compute its eigenvalues and
orthonormal eigenvectors. Verify Theorem~\ref{ch42:thm-courant} by evaluating
$R_{\mathbf{M}}(x)$ at $x=(1,0)$, $x=(1,1)$, $x=(1,-1)$ and checking each lies in
$[\lambda_1,\lambda_2]$, with equality exactly at the eigenvectors.
\end{exercise}

\begin{exercise}[Laplacian is PSD]
For the path graph on three vertices $1\!-\!2\!-\!3$, write $\Lap$ explicitly, exhibit it as
$\mathbf{B}^\top\mathbf{B}$ for the signed incidence matrix $\mathbf{B}$, and confirm
$x^\top\Lap x=(x_1-x_2)^2+(x_2-x_3)^2\ge 0$. Find $\ker\Lap$ and relate its dimension to the
number of connected components.
\end{exercise}

\begin{exercise}[Indicator trick]
A uniformly random permutation of $\{1,\dots,n\}$ has a \emph{fixed point} at $i$ if it maps
$i\mapsto i$. Using indicator variables and linearity, show the expected number of fixed
points is exactly $1$, independent of $n$. Then use Chebyshev to bound the probability of
seeing $\ge t$ fixed points.
\end{exercise}

\begin{exercise}[Boosting by repetition]
A Monte Carlo algorithm answers correctly with probability $2/3$ on every input. You run it
$2k+1$ independent times and take the majority vote. Using the Chernoff bound of
Theorem~\ref{ch42:thm-concentration}, show the majority is wrong with probability at most
$e^{-\Om(k)}$, and deduce how many repetitions guarantee error below $1/n$.
\end{exercise}

\begin{exercise}[Karger survival, warm-up]
In a graph with min cut of size $c$, every vertex has degree $\ge c$, so $m\ge nc/2$. Show
that a uniformly random edge avoids a fixed min cut with probability $\ge 1-\frac{2}{n}$.
Chain these across $n-2$ contractions to derive the $\ge 2/(n(n-1))$ survival probability
quoted in Section~\ref{ch42:probability}.
\end{exercise}
